\begin{table}[H]
  \centering
  \caption{Sensibilidad del núcleo SHAP y de su dirección a los umbrales
  elegidos (Sección~\ref{subsec:importancia_resultados}). ``(base)'': criterio
  usado en el texto. El núcleo de 22 variables se conserva por completo con
  criterios más laxos; con los más estrictos queda un subconjunto.}
  \label{tab:sensibilidad_shap}
  \footnotesize
  \begin{tabular}{ccccc}
    \toprule
    \multicolumn{5}{l}{\textbf{A. Núcleo: masa SHAP acumulada y mínimo de combinaciones (de 10)}} \\
    \midrule
    Masa & Mín. comb. & Núcleo & Conserva de las 22 & Universales en las 4 fuentes (de 7) \\
    \midrule
    70\% & 4 & 19 & 17 & 6 \\
    70\% & 5 & 13 & 13 & 4 \\
    70\% & 6 & 11 & 11 & 3 \\
    80\% & 4 & 30 & 22 & 7 \\
    80\% & 5 & 22 (base) & 22 & 7 \\
    80\% & 6 & 20 & 20 & 5 \\
    90\% & 4 & 50 & 22 & 7 \\
    90\% & 5 & 35 & 22 & 7 \\
    90\% & 6 & 31 & 22 & 7 \\
    \midrule
    \multicolumn{5}{l}{\textbf{B. Dirección: |$\rho$| mínimo y proporción mínima de signo}} \\
    \midrule
    $|\rho|$ mín. & Prop. mín. & Consistentes (de 20) & & \\
    \midrule
    0.05 & 0.70 & 12 & & \\
    0.05 & 0.80 & 10 & & \\
    0.05 & 0.90 & 7 & & \\
    0.10 & 0.70 & 12 & & \\
    0.10 & 0.80 & 10 (base) & & \\
    0.10 & 0.90 & 7 & & \\
    0.15 & 0.70 & 12 & & \\
    0.15 & 0.80 & 10 & & \\
    0.15 & 0.90 & 6 & & \\
    0.20 & 0.70 & 12 & & \\
    0.20 & 0.80 & 10 & & \\
    0.20 & 0.90 & 6 & & \\
    \bottomrule
  \end{tabular}